
\section{\BookRecordIndex{T0139}格罗莫夫紧性}\label{sec:5-3}
\zhindex{geluomofujinxing}{格罗莫夫紧性}

本节将证明，每一列能量一致有界的 $J$-全纯球面都有\BookRecordIndex{T0141}格罗莫夫收敛子列。

\begin{theorem}[\BookRecordIndex{T0139}格罗莫夫紧性]\label{thm:5-3-1}
设 $(M,\omega)$ 是紧\BookRecordIndex{T0358}辛流形，且
$J^\nu\in\mathcal J_\tau(M,\omega)$ 是 $M$ 上一列被 $\omega$ \BookRecordIndex{T0372}驯服的\BookRecordIndex{T0088}殆复结构，
它在 $C^\infty$ 拓扑下收敛到 $J\in\mathcal J_\tau(M,\omega)$。设
$u^\nu:S^2\to M$ 是一列 $J^\nu$-全纯球面，满足
$\sup_\nu E(u^\nu)<\infty$，并设
$\mathbf z^\nu=(z_1^\nu,\ldots,z_n^\nu)$ 是 $S^2$ 中两两不同的点所成的一列
$n$-元组。则 $(u^\nu,\mathbf z^\nu)$ 有\BookRecordIndex{T0141}格罗莫夫收敛子列。
\end{theorem}

由定理~\ref{thm:5-3-1} 和定理~\ref{thm:5-2-2} 可知，对每个与 $\omega$ \BookRecordIndex{T0348}相容的
\BookRecordIndex{T0088}殆复结构 $J$，在辛面积有固定上界时，只有有限多个同伦类能够由 $J$-全纯球面表示。
（命题~\ref{prop:4-1-5} 给出了一个更容易的证明。）因此，由习题~\ref{ex:5-1-2}，
只有有限多个这样的同伦类能够由亏格为零的\BookRecordIndex{T0340}稳定映射表示。下面的推论把这一观察推广到
\BookRecordIndex{T0088}殆复结构族。

\begin{corollary}\label{cor:5-3-2}
设 $\Lambda$ 是紧度量空间，且映射
\[
  \Lambda\longrightarrow\mathcal J(M,\omega):\lambda\longmapsto J_\lambda
\]
相对于 $\mathcal J(M,\omega)$ 上的 $C^\infty$ 拓扑连续。则对每个 $c>0$，只有有限多个
满足 $\langle[\omega],A\rangle\le c$ 的同伦类 $A\in\pi_2(M)$，能够对某个
$\lambda\in\Lambda$ 由 $J_\lambda$-全纯\BookRecordIndex{T0340}稳定映射（亏格为零）表示。
\end{corollary}

\begin{proof}
由命题~\ref{prop:4-1-4}，只需证明：只有有限多个满足
$\langle[\omega],A\rangle\le c$ 的同伦类 $A\in\pi_2(M)$，能够对某个 $\lambda$ 由
$J_\lambda$-全纯球面表示。反设不然，则存在序列 $\lambda^\nu\in\Lambda$，以及一列
$J_{\lambda^\nu}$-全纯球面 $u^\nu:S^2\to M$；其中每个球面表示一个不同的同伦类
$A^\nu$，且 $\langle[\omega],A^\nu\rangle\le c$。由定理~\ref{thm:5-3-1}，
$u^\nu$ 有\BookRecordIndex{T0141}格罗莫夫收敛子列 $u^{\nu_j}$。由定理~\ref{thm:5-2-2}，当 $j$ 充分大时，
该子列的各项表示同一个同伦类。这与假设矛盾。
\end{proof}

\begin{exercise}\label{ex:5-3-3}
用命题~\ref{prop:4-1-5} 的方法，给出推论~\ref{cor:5-3-2} 的另一个简单得多的证明；
即使映射
$\Lambda\to\mathcal J(M,\omega):\lambda\mapsto J_\lambda$
仅相对于 $\mathcal J(M,\omega)$ 上的 $C^0$ 拓扑连续，这个证明仍然成立。
\end{exercise}

我们先证明没有\BookRecordIndex{T0066}标记点时的\BookRecordIndex{T0192}紧性定理。这个论证大量依赖第~\ref{sec:4-7} 节的结果。

\begin{proof}[定理~\ref{thm:5-3-1} 在 $n=0$ 时的证明]
设 $u^\nu:S^2\to M$ 是一列 $J^\nu$-全纯球面。不失一般性，可以假设
$E(u^\nu)$ 收敛，并把其极限记为
\[
  E=\lim_{\nu\to\infty}E(u^\nu).
\]
为了实施归纳，把有 $N$ 个顶点的树 $T$ 写成向量
$\mathbf j=(j_2,\ldots,j_N)$ 会很方便，其中对每个 $i$ 都有 $j_i<i$，如
习题~\ref{ex:5-1-3} 所述。我们将用归纳法构造：
\begin{enumerate}[label=(\alph*),leftmargin=3em]
\item 一个 $(3N-2)$-元组
\[
  \mathbf u=(u_1,\ldots,u_N;j_2,\ldots,j_N;z_2,\ldots,z_N),
\]
其中 $u_i:S^2\to M$ 是 $J$-全纯球面；当 $i\ge2$ 时，$j_i<i$ 是正整数；
$z_i\in\mathbb C$ 且 $|z_i|\le1$；$u_1$ 非常值；当 $k\ne k'$ 且
$j_k=j_{k'}$ 时，就有 $z_k\ne z_{k'}$
\sourcefootnote{原书只写“只要 $j_k=j_{k'}$，就有 $z_k\ne z_{k'}$”，遗漏了
$k\ne k'$；按字面取 $k=k'$ 会自相矛盾，而这里所需且由后文唯一确定的条件，是同一
母顶点上的不同子分支具有互异的附着点，故补上 $k\ne k'$。}；并且当
$i=2,\ldots,N$ 时，
$u_{j_i}(z_i)=u_i(\infty)$；
\item 对 $i=1,\ldots,N$，有限子集 $Z_i\subset B_1$，其中
$Z_1\subset\{0\}$，并且只要 $u_i$ 是常值映射，就有 $\#Z_i\ge2$；
\item 对 $i=1,\ldots,N$，\BookRecordIndex{T0040}Möbius 变换列
$\{\phi_i^\nu\}_{\nu\in\mathbb N}$；
\end{enumerate}
使某个适当的子列满足以下条件：
\begin{enumerate}[label=(\roman*),leftmargin=3em]
\item $u^\nu\circ\phi_1^\nu$ 在 $S^2\setminus Z_1$ 上\BookRecordIndex{T0416}逐紧一致收敛到 $u_1$。
当 $i=2,\ldots,N$ 时，$u^\nu\circ\phi_i^\nu$ 在 $\mathbb C\setminus Z_i$ 上
\BookRecordIndex{T0416}逐紧一致收敛到 $u_i$。

\item 对每个 $z\in Z_i$，极限
\[
  m_i(z)=\lim_{\varepsilon\to0}\lim_{\nu\to\infty}
  E\bigl(u^\nu\circ\phi_i^\nu;B_\varepsilon(z)\bigr)
\]
存在且为正。若 $Z_1=\{0\}$，则 $E=E(u_1)+m_1(0)$。若 $i\ge2$，则
$z_i\in Z_{j_i}$，且
\[
  m_{j_i}(z_i)=E(u_i)+\sum_{z\in Z_i}m_i(z).
\]

\item 当 $i=2,\ldots,N$ 时，$(\phi_{j_i}^\nu)^{-1}\circ\phi_i^\nu$ 在
$\mathbb C=S^2\setminus\{\infty\}$ 上\BookRecordIndex{T0416}逐紧一致收敛到 $z_i$。

\item $Z_N=\varnothing$，且
\[
  Z_i=\{z_k\mid i<k\le N,\ j_k=i\},
  \qquad i=1,\ldots,N-1.
\]
\end{enumerate}

解释一下这些条件。条件 (i) 和 (ii) 的第一部分合起来说明，每列
$u^\nu\circ\phi_i^\nu$ 恰好在 $Z_i$ 的各点\BookRecordIndex{T0243}冒泡。因此，第一个\BookRecordIndex{T0275}气泡 $u_1$ 至多有一个
\BookRecordIndex{T0243}冒泡点 $Z_1$，从而是树 $T$ 的端点；而由 (a)，当 $i>1$ 时，第 $i$ 个\BookRecordIndex{T0275}气泡 $u_i$
在点 $u_i(\infty)=u_{j_i}(z_i)$ 处附着于第 $j_i$ 个\BookRecordIndex{T0275}气泡。此外，(ii) 的第二部分
说明，$u_{j_i}$ 在 $z_i$ 处的全部能量，都由 $u_i$ 与它在 $Z_i$ 处的\BookRecordIndex{T0275}气泡之并捕获。
由于只有当 $i>j$ 时才会有 $j_i=j$，树 $T$ 不含圈。我们先构造至多有一个\BookRecordIndex{T0243}冒泡点的
$u_1$，再用归纳法构造 $u_i$ 和 $Z_i$，使 (i)--(iii) 成立。归纳完成时，(iv) 也将
成立。换言之，对每个 $i$，$u_i$ 的全部\BookRecordIndex{T0243}冒泡都将由满足 $k>i$ 且 $j_k=i$ 的分支
$u_k$ 捕获。注意，(b) 包含了稳定性条件：我们把 $u_1$ 构造成非常值映射；当 $u_i$
为常值映射时，我们要求 $\#Z_i\ge2$ 而不是 $3$，因为当 $i\ge2$
时，$\infty$ 是第 $i$ 个球面上的另一个\BookRecordIndex{T0191}节点
\sourcefootnote{原书此处写作 $i\ge1$；但根分支 $i=1$ 没有指向母分支的\BookRecordIndex{T0191}节点
$\infty$，而且前文已经规定 $u_1$ 非常值，故这里唯一相关的指标范围是 $i\ge2$。}。

在整个证明中，我们把 $u^\nu$ 看成从 $\mathbb C$ 到 $M$ 的函数，并用
$|du^\nu(z)|$ 表示微分
$du^\nu(z):\mathbb C\to T_{u^\nu(z)}M$ 关于 $\mathbb C$ 上通常度量的范数。于是，
对每个 $\nu$，$du^\nu(z)$ 关于 Fubini--Study 度量的范数有界，并且
\[
  |du^\nu(z)|_{\mathrm{FS}}=|du^\nu(z)|(1+|z|^2)
\]
（见习题~\ref{ex:4-2-3}）。因此，当 $|z|\to\infty$ 时，每个函数 $|du^\nu(z)|$ 都必须趋于
零。所以存在序列 $z^\nu\in\mathbb C$，使得
\[
  |du^\nu(z^\nu)|=\sup_{z\in\mathbb C}|du^\nu(z)|=:c^\nu.
\]
现在注意，序列
\[
  v^\nu(z):=u^\nu(z^\nu+z/c^\nu)
\]
满足 $\sup|dv^\nu|=1=|dv^\nu(0)|$，因而有一个子列在
$\mathbb C=S^2\setminus\{\infty\}$ 上\BookRecordIndex{T0416}逐紧一致收敛到一条非常值 \BookRecordIndex{T0003}$J$-全纯曲线 $v$。
再取子列，可以假设 $E(v^\nu;\mathbb C\setminus B_1)$ 收敛。于是，函数
\[
  u_1(z)=v(1/z),
  \qquad
  \phi_1^\nu(z)=z^\nu+\frac1{c^\nu z}
\]
以 $Z_1=\varnothing$ 或 $Z_1=\{0\}$ 满足 (i) 和 (ii)。当 $N=1$ 时，条件 (iii)
为空。若 $Z_1=\varnothing$，则 $u^\nu\circ\phi_1^\nu$ 在整个 $S^2$ 上一致且连同
所有阶导数一起收敛到 $u_1$，定理已经得证。因此，以下假设 $Z_1=\{0\}$。

设 $\ell\ge1$，并归纳假设：当 $i\le\ell$ 时，已经构造了
$u_i,j_i,z_i,Z_i$ 和 $\{\phi_i^\nu\}_{\nu\in\mathbb N}$，使它们满足
(i)--(iii) 而尚未满足 (iv)，其中以 $\ell$ 代替 $N$。于是存在 $j\le\ell$，使得
$Z_j\ne\{z_i\mid j<i\le\ell,\ j_i=j\}$。记
\begin{equation}\label{eq:5-3-1}
  Z_{j;\ell}:=Z_j\setminus\{z_i\mid j<i\le\ell,\ j_i=j\},
\end{equation}
任选 $z_{\ell+1}\in Z_{j;\ell}$，并把命题~\ref{prop:4-7-1} 应用于序列
$u^\nu\circ\phi_j^\nu$ 和点 $z_0=z_{\ell+1}$。仍照常以 $\nu$ 标记所得子列，便得到
满足命题~\ref{prop:4-7-1} 要求的元素 $\psi^\nu\in G$。特别地，
$u^\nu\circ\phi_j^\nu\circ\psi^\nu$ 在 $S^2\setminus Z$ 上\BookRecordIndex{T0416}逐紧一致收敛到一条
$J$-全纯球面 $u_{\ell+1}:S^2\to M$，其中 $Z\subset B_1$ 是有限集。
命题~\ref{prop:4-7-1} 还表明，$\psi^\nu$ 在
$\mathbb C=S^2\setminus\{\infty\}$ 上\BookRecordIndex{T0416}逐紧一致收敛到 $z_{\ell+1}$。此外，由
命题~\ref{prop:4-7-2}，
\[
  u_{\ell+1}(\infty)=u_j(z_{\ell+1}).
\]
这说明，当 $i\le\ell+1$ 时，以 $\ell+1$ 代替 $N$，并令
\[
  j_{\ell+1}:=j,
  \qquad Z_{\ell+1}:=Z,
  \qquad \phi_{\ell+1}^\nu:=\phi_j^\nu\circ\psi^\nu,
\]
条件 (i)--(iii) 均成立。归纳构造完成。

我们断言，这一归纳在有限步后终止。为此，在第 $\ell$ 步考察树
$T_\ell=(j_2,\ldots,j_\ell)$，赋予权重
\[
  m(j;\ell)=E(u_j)+\sum_{z\in Z_{j;\ell}}m_j(z),
  \qquad j=1,\ldots,\ell,
\]
其中 $Z_{j;\ell}$ 由式~\eqref{eq:5-3-1} 给出。（\BookRecordIndex{T0178}加权树定义于
习题~\ref{ex:5-1-2}。）可以把这个量理解为 $u_j$ 的能量，加上那些尚未被放大成
满足 $j<i\le\ell$ 的\BookRecordIndex{T0275}气泡 $u_i$ 的\BookRecordIndex{T0243}冒泡点能量。注意，只有当 $u_j$ 是常值映射、且它的
所有\BookRecordIndex{T0243}冒泡点都已处理完毕时，才有 $m(j;\ell)=0$。由 (b)，这说明 $j$ 是\BookRecordIndex{T0178}加权树
$T_\ell$ 的稳定分支。此外，$\hbar$ 是各非零权重 $m(j;\ell)$ 的统一正下界
\sourcefootnote{原书称 $\hbar$ 为这些非零权重的“最小值”；命题~\ref{prop:4-1-4}
只保证每个非零权重至少为 $\hbar$，下文也只需要这一统一正下界，故据此改正。}。
我们接着断言
\begin{equation}\label{eq:5-3-2}
  \sum_{j=1}^{\ell}m(j;\ell)=E.
\end{equation}
当 $\ell=1$ 时，这显然成立，因为
$E=E(u_1)+m_1(0)=m(1;1)$。现设式~\eqref{eq:5-3-2} 已对某个整数 $\ell\ge1$
成立。注意，只要 $j\ne j_{\ell+1}$ 且 $j\le\ell$，就有
$m(j;\ell+1)=m(j;\ell)$，而
\[
  m(j_{\ell+1};\ell+1)
  =m(j_{\ell+1};\ell)-m_{j_{\ell+1}}(z_{\ell+1}),
\]
并且由取 $i=\ell+1$ 的 (ii)，
\[
  m(\ell+1;\ell+1)=m_{j_{\ell+1}}(z_{\ell+1}).
\]
所以，以 $\ell+1$ 代替 $\ell$ 时，式~\eqref{eq:5-3-2} 仍成立。这样就对每个整数
$\ell\ge1$ 证明了式~\eqref{eq:5-3-2}。由习题~\ref{ex:5-1-2}，$T_\ell$ 的边数
不超过 $2E/\hbar$，故归纳必定终止。

于是，我们已经证明存在 $i=1,\ldots,N$ 时的
$u_i,j_i,z_i$ 和 $\{\phi_i^\nu\}_{\nu\in\mathbb N}$，使某个适当的子列满足
(i)--(iv)。在这些假设下，我们断言子列 $u^\nu$ \BookRecordIndex{T0141}格罗莫夫收敛到\BookRecordIndex{T0340}稳定映射
$\mathbf u$。先验证 $\mathbf u$ 稳定。注意，对顶点 $j$ 和 $i>j$，它们由一条边连接
当且仅当 $j=j_i$。此外，用定义~\ref{def:5-1-1} 的记号，
$z_{j_i i}=z_i$ 且 $z_{i j_i}=\infty$。因此，对顶点
$j\in\{2,\ldots,N\}$，其奇异集是 $Z_j\cup\{\infty\}$，其中 $Z_j$ 如 (iv) 所示；
而顶点 $1$ 的奇异集是 $Z_1$。若 $u_i$ 是常值映射，则 $i\ge2$ 且 $\#Z_i\ge2$。
所以 $\mathbf u$ 如所断言是稳定的。

还须验证定义~\ref{def:5-2-1} 的四个条件。当 $n=0$ 时，（\BookRecordIndex{T0066}标记点）公理为空；
（映射）和（重标度）公理由 (i) 和 (a) 得出。因此，只须验证（能量）公理。

用 $T_i$ 表示如下顶点 $i'$ 的集合：从 $j_i$ 出发，沿着以 $j_i$ 到 $i$ 的边为首边的
\BookRecordIndex{T0064}边链可以到达 $i'$。在第~\ref{sec:5-1} 节的记号中，$T_i=T_{j_i i}$。这些集合可以
按 $i=N,N-1,\ldots,1$ 递归定义为
\[
  T_i:=\{i\}\cup\bigcup_{j_k=i}T_k.
\]
特别地，若不存在满足 $k>i$ 且 $j_k=i$ 的 $k$，则 $T_i=\{i\}$。先对 $T_i$ 的元素
个数作归纳，证明
\begin{equation}\label{eq:5-3-3}
  m_{j_i}(z_i)=\sum_{i'\in T_i}E(u_{i'}).
\end{equation}
当 $Z_i=\varnothing$，因而 $T_i=\{i\}$ 时，这立即来自 (ii)。因此，假设这个公式已经
对 $\#T_i\le m$ 得证。设 $\#T_i=m+1$，并令
\[
  \{i_1,\ldots,i_s\}:=\{k>i\mid j_k=i\}.
\]
于是 $\{i_1,\ldots,i_s\}\cup\{j_i\}$ 是与 $i$ 相邻的顶点集，而且对
$r=1,\ldots,s$，有 $\#T_{i_r}\le m$。所以，由归纳假设，
\[
  m_i(z_{i_r})=\sum_{i'\in T_{i_r}}E(u_{i'}),
  \qquad r=1,\ldots,s.
\]
另一方面，由 (ii)，
\[
  m_{j_i}(z_i)=E(u_i)+\sum_{r=1}^s m_i(z_{i_r}).
\]
代入前面的公式，并使用
$T_i=\{i\}\cup\bigcup_rT_{i_r}$，便得到式~\eqref{eq:5-3-3}。它等价于
式~\eqref{eq:5-2-1} 取 $\alpha=j_i$、$\beta=i$ 的情形。

现在考察 $\alpha=i$、$\beta=j_i$ 的情形。此时
$T_{\alpha\beta}=\{1,\ldots,N\}\setminus T_i$ 且 $z_{\alpha\beta}=\infty$，
所以式~\eqref{eq:5-2-1} 成为
\begin{equation}\label{eq:5-3-4}
  m_i(\infty):=\lim_{R\to\infty}\lim_{\nu\to\infty}
  E\bigl(u^\nu\circ\phi_i^\nu;\mathbb C\setminus B_R\bigr)
  =\sum_{i'\notin T_i}E(u_{i'}).
\end{equation}
为证明这一点，注意
\[
  \begin{aligned}
  E-m_i(\infty)
  &=\lim_{R\to\infty}\lim_{\nu\to\infty}
    E\bigl(u^\nu\circ\phi_i^\nu;B_R\bigr)\\
  &=E(u_i)+\sum_{z\in Z_i}m_i(z)\\
  &=m_{j_i}(z_i)\\
  &=\sum_{i'\in T_i}E(u_{i'}).
  \end{aligned}
\]
第一个等号来自 $E:=\lim_{\nu\to\infty}E(u^\nu)$ 和 $m_i(\infty)$ 的定义；
第二个来自 $m_i(z)$ 的定义；第三个来自 (ii)；第四个来自式~\eqref{eq:5-3-3}。
此外，在式~\eqref{eq:5-3-2} 中取 $\ell=N$，得到
$E=\sum_{i=1}^NE(u_i)$；把它与上式结合，便得到式~\eqref{eq:5-3-4}。
定理~\ref{thm:5-3-1} 在 $n=0$ 时得证。
\end{proof}

下面考虑有\BookRecordIndex{T0066}标记点时的情形。归纳假设序列
$(u^\nu,z_1^\nu,\ldots,z_k^\nu)$ \BookRecordIndex{T0141}格罗莫夫收敛到\BookRecordIndex{T0340}稳定映射
\[
  (\mathbf u,\mathbf z)
  =\bigl(\{u_\alpha\}_{\alpha\in T},
    \{z_{\alpha\beta}\}_{\alpha E\beta},
    \{z_i\}_{1\le i\le k}\bigr)
  \in\mathcal{SC}_{0,k}(M;J)
\]
其中 $k\ge0$；对每个 $\alpha\in T$，对应的变换列为 $\phi_\alpha^\nu\in G$。
我们来考察再加入一列\BookRecordIndex{T0066}标记点
$\zeta^\nu\in S^2\setminus\{z_1^\nu,\ldots,z_k^\nu\}$
\sourcefootnote{原书在此处及随后引理~\ref{lem:5-3-4} 的假设中，两次把集合下标的
起点印作未定义的 $i$；此时已有且新增点必须避开的是全部 $k$ 个\BookRecordIndex{T0066}标记点，故统一改为
$1$。} 后发生的情形。取子列后，
可以假设对所有 $\alpha\in T$，极限
\begin{equation}\label{eq:5-3-5}
  \zeta_\alpha=\lim_{\nu\to\infty}(\phi_\alpha^\nu)^{-1}(\zeta^\nu)
\end{equation}
都存在。回忆
$Z_\alpha=\{z_{\alpha\beta}\mid\beta\in T,\ \alpha E\beta\}$，它是
$\alpha$-球面上相邻\BookRecordIndex{T0275}气泡的附着点集。

\begin{lemma}\label{lem:5-3-4}
设 $(u^\nu,z_1^\nu,\ldots,z_k^\nu)$ 通过 $\phi_\alpha^\nu\in G$ \BookRecordIndex{T0141}格罗莫夫收敛到
\BookRecordIndex{T0340}稳定映射 $(\mathbf u,\mathbf z)\in\mathcal{SC}_{0,k}(M;J)$，并且序列
$\zeta^\nu\in S^2\setminus\{z_1^\nu,\ldots,z_k^\nu\}$ 满足
式~\eqref{eq:5-3-5}。则以下条件恰有一个成立：
\begin{enumerate}[label=(\Roman*),leftmargin=3em]
\item 存在唯一的顶点 $\alpha\in T$，使得
\[
  \zeta_\alpha\notin
  Y_\alpha(\mathbf u,\mathbf z)
  =Z_\alpha\cup\{z_i\mid\alpha_i=\alpha\}.
\]

\item 存在唯一的指标 $i\in\{1,\ldots,k\}$，使得
$\zeta_{\alpha_i}=z_i$。

\item 存在唯一的边 $\alpha E\beta$，使得
$\zeta_\alpha=z_{\alpha\beta}$ 且 $\zeta_\beta=z_{\beta\alpha}$。
\end{enumerate}
\end{lemma}

\begin{proof}
先注意，集合族 $\{\zeta_\alpha\}_{\alpha\in T}$ 满足
\begin{equation}\label{eq:5-3-6}
  \alpha E\gamma,\quad \zeta_\alpha\ne z_{\alpha\gamma}
  \quad\Longrightarrow\quad
  \zeta_\gamma=z_{\gamma\alpha}.
\end{equation}
这是式~\eqref{eq:5-3-5} 以及下述事实的推论：序列
$\phi_{\gamma\alpha}^\nu=(\phi_\gamma^\nu)^{-1}\circ\phi_\alpha^\nu$
在 $S^2\setminus\{z_{\alpha\gamma}\}$ 上\BookRecordIndex{T0416}逐紧一致收敛到 $z_{\gamma\alpha}$。
现设存在某个 $\alpha\in T$，使得只要 $\alpha E\gamma$，就有
$\zeta_\alpha\ne z_{\alpha\gamma}$。于是式~\eqref{eq:5-3-6} 表明，对每个
$\beta\ne\alpha$ 都有 $\zeta_\beta\in Z_\beta$。事实上，给定
$\beta\ne\alpha$，选取从 $\gamma_0=\alpha$ 到 $\gamma_m=\beta$ 的\BookRecordIndex{T0064}边链
$\gamma_0,\ldots,\gamma_m\in T$。用式~\eqref{eq:5-3-6} 和
$\zeta_{\gamma_0}\ne z_{\gamma_0\gamma_1}$ 作归纳，得到
\[
  \zeta_{\gamma_i}=z_{\gamma_i\gamma_{i-1}},
  \qquad i=1,\ldots,m.
\]
这同时证明了 (I) 和 (II) 中的唯一性。

再证明 (III) 中的唯一性。设有两对 $\alpha E\beta$ 和
$\alpha' E\beta'$，满足
\[
  \zeta_\alpha=z_{\alpha\beta},\qquad
  \zeta_\beta=z_{\beta\alpha},\qquad
  \zeta_{\alpha'}=z_{\alpha'\beta'},\qquad
  \zeta_{\beta'}=z_{\beta'\alpha'}.
\]
令 $\gamma_0,\ldots,\gamma_m\in T$ 是从 $\gamma_0=\alpha$ 到
$\gamma_m=\beta'$ 的\BookRecordIndex{T0064}边链。必要时交换 $\alpha$ 和 $\beta$，并在\BookRecordIndex{T0064}边链中添加一项，
可以假设 $\gamma_1=\beta$；同理，可以假设 $\gamma_{m-1}=\alpha'$。若 $m\ge2$，
则 $\zeta_{\gamma_1}=z_{\gamma_1\gamma_0}\ne z_{\gamma_1\gamma_2}$，由
式~\eqref{eq:5-3-6}，$\zeta_{\gamma_2}=z_{\gamma_2\gamma_1}$。归纳可得
$\zeta_{\gamma_i}=z_{\gamma_i\gamma_{i-1}}$ 对所有 $i\ge1$ 成立，因而
\[
  \zeta_{\alpha'}=z_{\gamma_{m-1}\gamma_{m-2}}
  \ne z_{\alpha'\beta'}.
\]
这与假设矛盾，故 $m=1$，从而
$\{\alpha,\beta\}=\{\alpha',\beta'\}$。

现在假设 (I) 和 (II) 都不成立。则对每个 $\alpha\in T$，存在
$\beta\in T$，使 $\alpha E\beta$ 且 $\zeta_\alpha=z_{\alpha\beta}$。我们证明存在一对
$\alpha E\beta$，同时满足
$\zeta_\alpha=z_{\alpha\beta}$ 和 $\zeta_\beta=z_{\beta\alpha}$。任取一对
$\alpha_0E\alpha_1$，使 $\zeta_{\alpha_0}=z_{\alpha_0\alpha_1}$。若
$\zeta_{\alpha_1}\ne z_{\alpha_1\alpha_0}$，则
$\zeta_{\alpha_1}=z_{\alpha_1\alpha_2}$ 对某个 $\alpha_2\ne\alpha_0$ 成立；由此可以构造
序列 $\alpha_i\in T$，满足 $\alpha_iE\alpha_{i+1}$ 且
$\zeta_{\alpha_i}=z_{\alpha_i\alpha_{i+1}}$。这个序列只有在某个 $i$ 满足
$\zeta_{\alpha_{i+1}}=z_{\alpha_{i+1}\alpha_i}$
\sourcefootnote{原书此处印作 $\zeta_{i+1}$；但式~\eqref{eq:5-3-5} 只定义了按树顶点
$\alpha$ 标记的量 $\zeta_\alpha$，故补为 $\zeta_{\alpha_{i+1}}$。}
时才会终止，而因为 $T$ 是有限集，它必定终止。因此，(I)、(II)、(III) 中至少有一个
成立。容易看出，这三个条件两两互斥。
\end{proof}

还须说明怎样在上述三种情形中容纳新增的\BookRecordIndex{T0066}标记点。情形 (I) 很直接；但在另外两种
情形中，必须给树添加一个新顶点。

\begin{proof}[定理~\ref{thm:5-3-1} 在 $n>0$ 时的证明]
设 $(u^\nu,z_1^\nu,\ldots,z_n^\nu)$ 是一列 $J^\nu$-全纯球面，每个球面都有
$n$ 个两两不同的\BookRecordIndex{T0066}标记点。我们对 $k$ 作归纳，证明
$(u^\nu,z_1^\nu,\ldots,z_k^\nu)$ 的某个子列\BookRecordIndex{T0141}格罗莫夫收敛到\BookRecordIndex{T0340}稳定映射
\[
  (\mathbf u,\mathbf z)
  =\bigl(\{u_\alpha\}_{\alpha\in T},
    \{z_{\alpha\beta}\}_{\alpha E\beta},
    \{z_i\}_{1\le i\le k}\bigr)
  \in\mathcal{SC}_{0,k}(M;J),
\]
其对应的 \BookRecordIndex{T0040}Möbius 变换为 $\phi_\alpha^\nu$。$k=0$ 的情形已经证明。设 $k\ge1$，
并归纳假设以 $k-1$ 代替 $k$ 时结论成立。必要时再取子列，可以假设对所有 $\alpha$，
极限
\begin{equation}\label{eq:5-3-7}
  z_{\alpha k}:=\lim_{\nu\to\infty}(\phi_\alpha^\nu)^{-1}(z_k^\nu)
\end{equation}
都存在，因而可以应用引理~\ref{lem:5-3-4}。所以，条件 (I)--(III) 恰有一个成立。

若 (I) 成立，则存在唯一元素 $\alpha_k\in T$，使得
$z_{\alpha_k k}\notin Y_{\alpha_k}(\mathbf u,\mathbf z)$。此时在
$\alpha_k$-球面上引入新\BookRecordIndex{T0066}标记点
\[
  z_k^{\mathrm{new}}:=z_{\alpha_k k}.
\]

若 (II) 成立，则存在唯一指标 $i\le k-1$，使得 $z_{\alpha_i k}=z_i$。此时选取
\BookRecordIndex{T0040}Möbius 变换列 $\psi^\nu\in G$，满足
\[
  \psi^\nu(0)=z_i^\nu,
  \qquad
  \psi^\nu(1)=z_k^\nu,
  \qquad
  \psi^\nu(\infty)=\phi_{\alpha_i}^\nu(w),
\]
其中选取 $w\in S^2\setminus\{z_i\}$，使得对所有 $\nu$ 都有
$\phi_{\alpha_i}^\nu(w)\notin\{z_i^\nu,z_k^\nu\}$。于是序列
$(\phi_{\alpha_i}^\nu)^{-1}\circ\psi^\nu(0)$ 和
$(\phi_{\alpha_i}^\nu)^{-1}\circ\psi^\nu(1)$ 都收敛到 $z_i$，而序列
$(\psi^\nu)^{-1}\circ\phi_{\alpha_i}^\nu(w)$ 恒为 $\infty$。所以，由引理 D.1.4，
$(\phi_{\alpha_i}^\nu)^{-1}\circ\psi^\nu$ 在
$\mathbb C=S^2\setminus\{\infty\}$ 上\BookRecordIndex{T0416}逐紧一致收敛到 $z_i$。

现在给树 $T$ 添加一个新顶点 $\gamma$；它只与 $\alpha_i$ 相连，并承载\BookRecordIndex{T0066}标记点
$z_i$ 和 $z_k$。令相应的变换列 $\phi_\gamma^\nu$ 等于 $\psi^\nu$。于是
\[
  (\phi_\gamma^\nu)^{-1}(z_i^\nu)=0,
  \qquad
  (\phi_\gamma^\nu)^{-1}(z_k^\nu)=1,
\]
并且 $(\phi_{\alpha_i}^\nu)^{-1}\circ\phi_\gamma^\nu$ 在
$S^2\setminus\{\infty\}$ 上\BookRecordIndex{T0416}逐紧一致收敛到 $z_i$。由于
$z_i\notin Z_{\alpha_i}$，序列 $u^\nu\circ\phi_\gamma^\nu$ 在
$\mathbb C=S^2\setminus\{\infty\}$ 上\BookRecordIndex{T0416}逐紧一致收敛到常值函数
$u_\gamma(z)\equiv u_{\alpha_i}(z_i)$。新的\BookRecordIndex{T0340}稳定映射由
$T^{\mathrm{new}}:=T\cup\{\gamma\}$ 以及
\[
  z_{\gamma\alpha_i}^{\mathrm{new}}:=\infty,
  \qquad z_{\alpha_i\gamma}^{\mathrm{new}}:=z_i,
  \qquad \alpha_k^{\mathrm{new}}:=\gamma,
  \qquad z_k^{\mathrm{new}}:=1,
  \qquad \alpha_i^{\mathrm{new}}:=\gamma,
  \qquad z_i^{\mathrm{new}}:=0
\]
给出。这个新\BookRecordIndex{T0275}气泡是“\BookRecordIndex{T0384}幽灵气泡”，即对应映射 $u_\gamma$ 为常值。为验证
定义~\ref{def:5-2-1} 对新顶点 $\gamma$ 的（能量）公理，注意
$\gamma$-球面只有一个\BookRecordIndex{T0191}节点 $z_{\gamma\alpha_i}^{\mathrm{new}}=\infty$。由于极限
$u_\gamma$ 为常值，对每个 $R>0$，
$E(u^\nu\circ\phi_\gamma^\nu;B_R)$ 都收敛到零，因而
\[
  \lim_{\nu\to\infty}
  E\bigl(u^\nu\circ\phi_\gamma^\nu;\mathbb C\setminus B_R\bigr)
  =\lim_{\nu\to\infty}E(u^\nu)=E(\mathbf u).
\]
这证明了（能量）公理。

若 (III) 成立，则存在唯一的边 $\alpha E\beta$，使
$z_{\alpha k}=z_{\alpha\beta}$ 且 $z_{\beta k}=z_{\beta\alpha}$。我们证明存在
\BookRecordIndex{T0040}Möbius 变换列 $\psi^\nu\in G$，满足以下条件：
\begin{enumerate}[label=(\roman*),leftmargin=3em]
\item 对所有 $\nu$，$(\psi^\nu)^{-1}(z_k^\nu)=1$。

\item $(\phi_\alpha^\nu)^{-1}\circ\psi^\nu$ 在 $S^2\setminus\{\infty\}$ 上
\BookRecordIndex{T0416}逐紧一致收敛到 $z_{\alpha\beta}$，而
$(\phi_\beta^\nu)^{-1}\circ\psi^\nu$ 在 $S^2\setminus\{0\}$ 上\BookRecordIndex{T0416}逐紧一致收敛到
$z_{\beta\alpha}$。

\item $u^\nu\circ\psi^\nu$ 在 $S^2\setminus\{0,\infty\}$ 上\BookRecordIndex{T0416}逐紧一致收敛到
$u_\alpha(z_{\alpha\beta})=u_\beta(z_{\beta\alpha})$。

\item 对每个 $r>0$，
\[
  \lim_{\nu\to\infty}E(u^\nu\circ\psi^\nu;B_r)=m_{\alpha\beta}(\mathbf u),
  \qquad
  \lim_{\nu\to\infty}E(u^\nu\circ\psi^\nu;\mathbb C\setminus B_r)
  =m_{\beta\alpha}(\mathbf u).
\]
\end{enumerate}
为证明这一点，考察序列
\[
  \phi_{\alpha\beta}^\nu
  :=(\phi_\alpha^\nu)^{-1}\circ\phi_\beta^\nu,
  \qquad
  \xi^\nu:=(\phi_\beta^\nu)^{-1}(z_k^\nu),
\]
于是 $\xi^\nu\to z_{\beta\alpha}$ 且
$\phi_{\alpha\beta}^\nu(\xi^\nu)\to z_{\alpha\beta}$。不失一般性，可以假设
\[
  z_{\alpha\beta}=0,
  \qquad z_{\beta\alpha}=\infty.
\]
否则，选取 \BookRecordIndex{T0040}Möbius 变换 $\psi_\alpha$ 和 $\psi_\beta$，使
$\psi_\alpha(0)=z_{\alpha\beta}$、$\psi_\beta(\infty)=z_{\beta\alpha}$，并把
$\phi_\alpha^\nu$ 和 $\phi_\beta^\nu$ 分别替换为
$\phi_\alpha^\nu\circ\psi_\alpha$ 和 $\phi_\beta^\nu\circ\psi_\beta$。此外，不失
一般性，可以假设 $\phi_{\alpha\beta}^\nu(\infty)=\infty$。否则，由于
$\phi_{\beta\alpha}^\nu(\infty)\to\infty$，可以选取一列一致收敛到恒等映射的
等距变换 $\rho^\nu\in\operatorname{SO}(3)$，使得
$\rho^\nu\circ\phi_{\beta\alpha}^\nu(\infty)=\infty$，并把 $\phi_\beta^\nu$ 替换为
$\phi_\beta^\nu\circ(\rho^\nu)^{-1}$。经过这些调整，序列
$\phi_{\alpha\beta}^\nu$ 具有形式
\[
  \phi_{\alpha\beta}^\nu(z)=z^\nu+\varepsilon^\nu z,
\]
其中 $z^\nu,\varepsilon^\nu\in\mathbb C$ 都收敛到零，且对所有 $\nu$ 都有
$\varepsilon^\nu\ne0$。现在回忆
\[
  \lim_{\nu\to\infty}\xi^\nu=\infty,
  \qquad
  \lim_{\nu\to\infty}\varepsilon^\nu\xi^\nu=0.
\]
因此，对充分大的 $\nu$，重新定义序列 $\rho^\nu\in G$ 为
\[
  \rho^\nu(w):=\frac{w-z^\nu}{\varepsilon^\nu\xi^\nu}.
\]
它满足：
\begin{enumerate}[label=(\alph*),leftmargin=3em]
\item $\rho^\nu$ 在
$S^2\setminus\{z_{\alpha\beta}\}=S^2\setminus\{0\}$ 上\BookRecordIndex{T0416}逐紧一致收敛到 $\infty$；
\item $\rho^\nu\circ\phi_{\alpha\beta}^\nu$ 在
$S^2\setminus\{z_{\beta\alpha}\}=S^2\setminus\{\infty\}$ 上\BookRecordIndex{T0416}逐紧一致收敛到 $0$；
\item 对所有 $\nu$，
$\rho^\nu\circ(\phi_\alpha^\nu)^{-1}(z_k^\nu)=1$。
\end{enumerate}
令
\[
  \psi^\nu:=\phi_\alpha^\nu\circ(\rho^\nu)^{-1}.
\]
于是 $(\psi^\nu)^{-1}(z_k^\nu)=
\rho^\nu\circ(\phi_\alpha^\nu)^{-1}(z_k^\nu)=1$；由 (a)、(b)，
$(\phi_\alpha^\nu)^{-1}\circ\psi^\nu$ 在 $S^2\setminus\{\infty\}$ 上\BookRecordIndex{T0416}逐紧一致收敛到
$z_{\alpha\beta}$，而 $(\phi_\beta^\nu)^{-1}\circ\psi^\nu$ 在
$S^2\setminus\{0\}$ 上\BookRecordIndex{T0416}逐紧一致收敛到 $z_{\beta\alpha}$。所以，由
引理~\ref{lem:5-4-2}，$u^\nu\circ\psi^\nu$ 在 $S^2\setminus\{0,\infty\}$ 上
\BookRecordIndex{T0416}逐紧一致收敛到常值 $u_\alpha(z_{\alpha\beta})=u_\beta(z_{\beta\alpha})$，并且对每个
$r>0$，
\[
  \lim_{\nu\to\infty}E(u^\nu\circ\psi^\nu;B_r)=m_{\alpha\beta}(\mathbf u),
  \qquad
  \lim_{\nu\to\infty}E(u^\nu\circ\psi^\nu;\mathbb C\setminus B_r)
  =m_{\beta\alpha}(\mathbf u).
\]
所以序列 $\psi^\nu$ 如所断言满足 (i)--(iv)。

现在给树 $T$ 添加一个与重标度 $\phi_\gamma^\nu=\psi^\nu$ 对应的新顶点
$\gamma$。相应的极限映射 $u_\gamma$ 是常值映射：
\[
  u_\gamma(z)\equiv u_\alpha(z_{\alpha\beta})=u_\beta(z_{\beta\alpha}).
\]
这个\BookRecordIndex{T0275}气泡承载\BookRecordIndex{T0066}标记点 $z_k=1$。此外，对新的序列 $\phi_\gamma^\nu$，
式~\eqref{eq:5-3-7} 中的极限仍然存在。在新的\BookRecordIndex{T0340}稳定映射中，$\alpha$ 和 $\beta$ 不再相邻，
而由 $\gamma$ 隔开，并且
\[
  z_{\alpha\gamma}^{\mathrm{new}}=z_{\alpha\beta},
  \qquad z_{\beta\gamma}^{\mathrm{new}}=z_{\beta\alpha},
  \qquad z_{\gamma\alpha}^{\mathrm{new}}=\infty,
  \qquad z_{\gamma\beta}^{\mathrm{new}}=0,
  \qquad \alpha_k^{\mathrm{new}}=\gamma,
  \qquad z_k^{\mathrm{new}}=1.
\]
由 (i)--(iv)，定义~\ref{def:5-2-1} 的全部条件都成立。归纳和
定理~\ref{thm:5-3-1} 的证明完成。
\end{proof}
